% LaTeX テンプレート — 授業資料用 (LuaLaTeX を想定、両面: 奇数ページにヘッダー)
\documentclass[11pt,twoside,fleqn]{article}

% LuaLaTeX 向け設定: 日本語対応に luatexja-fontspec と fontspec を使用
\usepackage{fontspec}
\usepackage{luatexja-fontspec}

% 必要に応じて日本語フォントを明示（環境に応じて変更してください）
% 例: \setmainjfont{Noto Serif CJK JP}

% 単位表示の統一: siunitx を使用
\usepackage{siunitx}
\sisetup{detect-mode=true, per-mode=symbol}

% その他のパッケージ
\usepackage{tikz}
\usepackage{fancybox}
\usepackage{ascmac}
\usepackage{amsmath}
\usepackage{geometry}
\usepackage{fancyhdr}
\usepackage{lipsum} % サンプル用ダミーテキスト
\usepackage{ulem}
\usepackage{enumitem}
\usepackage{amssymb}
\usepackage{dashrule}

%Tikz
\usepackage{physmotion}

% 使い方: 以下を編集してください
\newcommand{\SubjectTitle}{重力による運動まとめ} %テーマ
\newcommand{\ClassRoomInfo}{\underline{\hspace{0.5cm}}年\underline{\hspace{0.5cm}}組　氏名\underline{\hspace{3cm}}}

% 余白設定
\geometry{top=20mm,bottom=25mm,left=20mm,right=20mm,heightrounded}

% ヘッダー設定: 奇数ページに表示（左: 科目と内容、右: 学年組欄）
\setlength{\headheight}{14pt}
\setlength{\headsep}{36pt}
\pagestyle{fancy}
\fancyhf{} % 既定のヘッダー/フッターをクリア
\fancyhead[LO]{\small 物理基礎「\SubjectTitle」} %Odd page Left headerを設定
\fancyhead[RO]{\small \ClassRoomInfo} % Odd page Right headerを設定
\renewcommand{\headrulewidth}{0.4pt}
\renewcommand{\footrulewidth}{0pt}
% フッター中央にページ番号を表示
\fancyfoot[C]{\small\thepage}

% 偶数ページはヘッダーを空にして、本文が上から始まるようにする
\fancyhead[LE,CE,RE]{}

% ドキュメント開始
\begin{document}
%復習：等加速度直線運動=====================================================================================
% \section{\textbf{復習}}
\begin{itembox}[l]{\textbf{〇等加速度直線運動}}
	\begin{minipage}[t]{0.70\linewidth}
	等加速度直線運動の変位・速度の式
	\Large
	\[
	\left(
	\begin{array}{ll}
		\text{⓪} \quad a = & \\
		\text{①} \quad v = & \\
		\text{②} \quad x = & \\
		\text{③} \quad v^2 - v_0^2 = 2ax &
	\end{array}
	\right.
	\]
	\begin{center}
		% \motiondiagram{初速度}{終速度}{加速度}：画像を挿入
		\motiondiagram{v_0}{v}{a}{t}
	\end{center}	
	\end{minipage}
	\hfill
	\begin{minipage}[t]{0.26\linewidth}
		\begin{shadebox}
			\begin{align*}
				v\,[\si{\metre\per\second}] &\text{：速度} \\
				v_0\,[\si{\metre\per\second}] &\text{：初速度} \\
				a\,[\si{\metre\per\second\squared}] &\text{：加速度} \\
				t\,[\si{\second}] &\text{：時間} \\
				x\,[\si{\metre}] &\text{：変位} \\
				g\,[\si{\metre\per\second\squared}] &\text{：重力加速度}\\
			\end{align*}
		\end{shadebox}
	\end{minipage}
\end{itembox}
%復習：鉛直投げ下ろし=====================================================================================
\begin{itembox}[l]{\textbf{〇鉛直投げ下ろし}}
	\begin{minipage}[t]{0.70\linewidth}
		鉛直投げ下ろし運動の変位・速度の式
		\Large
		\[
		\left(
		\begin{array}{ll}
			\text{⓪} \quad a = & \\
			\text{①} \quad v = & \\
			\text{②} \quad y = & \\
			\text{③} \quad v^2 - v_0^2 = 2gy &
		\end{array}
		\right.
		\]
		\normalsize
			\begin{shadebox}
				特徴：\\
				・加速度の向きと軸の向きが同じ\\
				\quad\to 加速度は$a=g$\\
				・等加速度直線運動の式について、$a=g$、$x=y$とすることで\\　鉛直投げ下ろし運動の式が得られる。\\
				・初速度$v_0=0$の時、自由落下になる。
			\end{shadebox}
	\end{minipage}
	\hfill
	\begin{minipage}[t]{0.26\linewidth}
		\begin{center}
			% \throwdowndiagram{初速度v_0}{終速度v}{加速度g}{時刻変数t}：画像を挿入
			\throwdowndiagram{v_0}{v}{g}{t}
		\end{center}
	\end{minipage}
\end{itembox}

%復習：鉛直投げ上げ=====================================================================================
\begin{itembox}[l]{\textbf{〇鉛直投げ上げ}}
	\begin{minipage}[t]{0.70\linewidth}
		鉛直投げ上げ運動の変位・速度の式
		\Large
		\begin{flalign*}
			&\left(
			\begin{array}{ll}
				\text{⓪} \quad a = & \\
				\text{①} \quad v = & \\
				\text{②} \quad y = & \\
				\text{③} \quad v^2 - v_0^2 = -2gy &
			\end{array}
			\right. &
		\end{flalign*}
		\normalsize
		\begin{shadebox}
				特徴：\\
				・加速度の向きと軸の向きが逆\\
				\quad\to 加速度は$a=-g$\\
				・等加速度直線運動の式について、$a=-g$、$x=y$とすることで\\　鉛直投げ上げ運動の式が得られる。\\
				・最高点では速度$v=0$になる。
		\end{shadebox}
	\end{minipage}
	\hfill
	\begin{minipage}[t]{0.26\linewidth}
		\begin{center}
			% \throwdowndiagram{初速度v_0}{終速度v}{加速度g}{時刻変数t}：画像を挿入
			\throwupdiagram{v_0}{v}{g}{t}
		\end{center}
	\end{minipage}
\end{itembox}
\clearpage
%復習：水平投射=====================================================================================
\begin{itembox}[l]{\textbf{〇水平投射}}
	\begin{minipage}[t]{0.60\linewidth}
		水平投射運動の変位・速度の式(水平方向)
		\Large
		\[
		\left(
		\begin{array}{ll}
			\text{⓪} \quad a = & \\
			\text{①} \quad v = & \\
			\text{②} \quad y = & 
		\end{array}
		\right.
		\]
		\normalsize
		水平投射運動の変位・速度の式(鉛直方向)
		\Large
		\[
		\left(
		\begin{array}{ll}
			\text{⓪} \quad a = & \\
			\text{①} \quad v = & \\
			\text{②} \quad y = &
		\end{array}
		\right.
		\]
		\normalsize
	\end{minipage}
	\hfill
	\begin{minipage}[t]{0.40\linewidth}
		\begin{center}
			% \throwdowndiagram{初速度v_0}{終速度v}{加速度g}{時刻変数t}：画像を挿入
			\projectilediagram{v_0}{v_x}{v_y}{v}{g}{t}
		\end{center}
	\end{minipage}
	\begin{shadebox}
		特徴：\\\\\\\\\\\\\\\\
	\end{shadebox}
\end{itembox}

\begin{itembox}[l]{\textbf{〇斜方投射}}
	\begin{minipage}[t]{0.50\linewidth}
		斜方投射運動の変位・速度の式(水平方向)
		\Large
		\[
		\left(
		\begin{array}{ll}
			\text{⓪} \quad a = & \\
			\text{①} \quad v = & \\
			\text{②} \quad y = & 
		\end{array}
		\right.
		\]
		\normalsize
		斜方投射運動の変位・速度の式(鉛直方向)
		\Large
		\[
		\left(
		\begin{array}{ll}
			\text{⓪} \quad a = & \\
			\text{①} \quad v = & \\
			\text{②} \quad y = &
		\end{array}
		\right.
		\]
		\normalsize
	\end{minipage}
	\hfill
	\begin{minipage}[t]{0.50\linewidth}
		\begin{center}
			% \throwdowndiagram{初速度v_0}{終速度v}{加速度g}{時刻変数t}：画像を挿入
			\obliqueprojectilediagram{v_0}{\theta}{v_{0x}}{v_{0y}}{g}{t}
		\end{center}
	\end{minipage}
	\begin{shadebox}
		特徴：\\\\\\\\\\\\\\\\
	\end{shadebox}
\end{itembox}



\clearpage
\end{document}

